% Figura: teorema del gradiente. (a) Trabajo de un gradiente entre a y b por
% dos caminos; (b) circulación nula en un circuito cerrado.
\newcommand{\curvasNivelElipt}{%
  \foreach \c/\pc in {0.6/100,1.2/80,1.8/60,2.4/40}{%
    \draw[azulprof!\pc!azulclaro, line width=0.8pt] (0,0) ellipse ({2*\c} and {\c});}
  \fill[azulprof] (0,0) circle (1.4pt);}
\begin{figure}[ht]
\centering
\begin{minipage}[t]{0.49\textwidth}\centering
\begin{tikzpicture}[>=Stealth, x=0.74cm, y=0.74cm, baseline=(current bounding box.north)]
  \clip (-5.1,-3.2) rectangle (5.1,3.2);
  \curvasNivelElipt
  % camino principal: x(t)=-4.2+6t, y(t)=-2.2+3.0t-1.4 sin(pi t)
  \draw[trayecto, domain=0:1, samples=60, smooth, variable=\t]
    plot ({-4.0+5.8*\t},{-2.0+3.0*\t-1.2*sin(180*\t)});
  % camino alternativo: segmento
  \draw[acento!60, densely dashed, line width=0.8pt] (-4.0,-2.0) -- (1.8,1.0);
  % vectores en puntos del camino principal
  \foreach \t in {0.18,0.52,0.86}{%
    \pgfmathsetmacro{\px}{-4.0+5.8*\t}\pgfmathsetmacro{\py}{-2.0+3.0*\t-1.2*sin(180*\t)}%
    \pgfmathsetmacro{\dx}{5.8}\pgfmathsetmacro{\dy}{3.0-1.2*pi*cos(180*\t)}%
    \pgfmathsetmacro{\dn}{sqrt(\dx*\dx+\dy*\dy)}%
    \pgfmathsetmacro{\gx}{-\px/2}\pgfmathsetmacro{\gy}{-2*\py}%
    \pgfmathsetmacro{\gn}{sqrt(\gx*\gx+\gy*\gy)}%
    \draw[->, azulprof, line width=1pt] (\px,\py) -- ++({0.8*\dx/\dn},{0.8*\dy/\dn});
    \draw[->, line width=0.9pt, cm2] (\px,\py) -- ++({1.2*\gx/\gn},{1.2*\gy/\gn});
    \fill (\px,\py) circle (0.8pt);}
  \fill[acento] (-4.0,-2.0) circle (1.2pt) node[below left=0pt, text=black] {$\vc{a}$};
  \fill[acento] (1.8,1.0) circle (1.2pt) node[above right=0pt, text=black] {$\vc{b}$};
  \node[font=\small, text=acento] at (-0.9,-2.1) {$C$};
  
  \node[fill=white, fill opacity=0.85, text opacity=1, inner sep=2pt, font=\footnotesize, anchor=south west]
       at (-5.0,2.45) {\textcolor{azulprof}{$\to\ \dl$}\quad\textcolor{cm2}{$\to\ \nab f$}};
\end{tikzpicture}\\[3pt]{\small\scshape (a)}\; $\displaystyle\int_C\nab f\cdot\dl=f(\vc{b})-f(\vc{a})$
\end{minipage}\hfill
\begin{minipage}[t]{0.49\textwidth}\centering
\begin{tikzpicture}[>=Stealth, x=0.74cm, y=0.74cm, baseline=(current bounding box.north)]
  \clip (-5.1,-3.2) rectangle (5.1,3.2);
  \curvasNivelElipt
  % circuito cerrado: elipse desplazada, sentido antihorario
  \def\cx{0.5}\def\cy{0.1}
  \draw[trayecto] ({\cx},{\cy}) ellipse (3.5 and 2.2);
  \foreach \s in {20,140,260}{%
    \pgfmathsetmacro{\px}{\cx+3.5*cos(\s)}\pgfmathsetmacro{\py}{\cy+2.2*sin(\s)}%
    \pgfmathsetmacro{\dx}{-3.5*sin(\s)}\pgfmathsetmacro{\dy}{2.2*cos(\s)}%
    \pgfmathsetmacro{\dn}{sqrt(\dx*\dx+\dy*\dy)}%
    \pgfmathsetmacro{\gx}{-\px/2}\pgfmathsetmacro{\gy}{-2*\py}%
    \pgfmathsetmacro{\gn}{sqrt(\gx*\gx+\gy*\gy)}%
    \draw[->, azulprof, line width=1pt] (\px,\py) -- ++({0.8*\dx/\dn},{0.8*\dy/\dn});
    \draw[->, line width=0.9pt, cm2] (\px,\py) -- ++({1.2*\gx/\gn},{1.2*\gy/\gn});
    \fill (\px,\py) circle (0.8pt);}
  \node[font=\small, text=acento] at (4.4,2.55) {$C$};
  \node[fill=white, fill opacity=0.85, text opacity=1, inner sep=2pt, font=\footnotesize, anchor=south west]
       at (-5.0,2.45) {\textcolor{azulprof}{$\to\ \dl$}\quad\textcolor{cm2}{$\to\ \nab f$}};
\end{tikzpicture}\\[3pt]{\small\scshape (b)}\; $\displaystyle\oint_C\nab f\cdot\dl=0$
\end{minipage}
\caption{Teorema del gradiente sobre el mapa de curvas de nivel de una
colina $f$ (ahora la cima es el máximo). (a)~Sobre la curva $C$ (y sobre
cualquier otra con los mismos extremos, p.\,ej. el segmento discontinuo) la
integral suma, en cada tramo, la proyección del gradiente sobre el desplazamiento
$\dl$: tramos cuesta arriba suman, tramos cuesta abajo restan, y el balance es
solo $f(\vc{b})-f(\vc{a})$. (b)~En un circuito cerrado, lo que se sube se baja:
la circulación es cero.}
\label{fig:teor_grad}
\end{figure}
